% Generated draft: requires mathematical and source review.
\documentclass[11pt]{article}
\usepackage[margin=25mm]{geometry}
\usepackage{fontspec}
\setmainfont{Latin Modern Roman}
\usepackage{amsmath,amssymb,mathtools,unicode-math}
\setmathfont{Latin Modern Math}
\usepackage{enumitem,needspace,xurl,hyperref}
\setlength{\parindent}{0pt}
\setlength{\parskip}{5pt}
\setlength{\emergencystretch}{3em}
\newcommand{\N}{\mathbb{N}}
\newcommand{\Z}{\mathbb{Z}}
\newcommand{\Q}{\mathbb{Q}}
\newcommand{\R}{\mathbb{R}}
\newcommand{\C}{\mathbb{C}}
\newcommand{\F}{\mathbb{F}}
\newcommand{\A}{\mathbb{A}}
\newcommand{\PP}{\mathbb{P}}
\newcommand{\E}{\mathbb{E}}
\newcommand{\eps}{\varepsilon}
\newcommand{\dd}{\,\mathrm{d}}
\begin{document}
\section*{UnsolvedMath ID: 800417}
\textbf{Source problem code:} catalogue entry
\subsection{Definitions and mathematical statement}
Let \(f:\PP^N\to\PP^N\) be a morphism of algebraic degree \(d\ge 2\), defined over \(\Q\), and let \(P\in\PP^N(\Q)\). Define its forward orbit by
\[
\mathcal O_f(P)=\{f^n(P):n\in\Z_{\ge0}\},\qquad f^0=\operatorname{id}.
\]
Assume that \(\overline{\mathcal O_f(P)}^{\mathrm{Zar}}=\PP^N\).

To make reduction precise, let \(G_f\) be the cofinite set of primes of good reduction for a fixed projective model of \(f\): at such a prime \(p\), reduction gives a morphism \(\widetilde f_p:\PP^N_{\F_p}\to\PP^N_{\F_p}\) of degree \(d\). Write \(\widetilde P_p\in\PP^N(\F_p)\) for the reduction of \(P\), and put
\[
L_p(f,P)=\#\{\widetilde f_p^{\,n}(\widetilde P_p):n\in\Z_{\ge0}\}.
\]
Thus \(L_p(f,P)\) counts all distinct points in the reduced forward orbit, including its initial nonperiodic part. For each real \(\epsilon>0\), define
\[
S_\epsilon(f,P)=\bigl\{p\in G_f:L_p(f,P)\le p^{(1/2+\epsilon)N}\bigr\}.
\]
\textbf{Question (supplied Question 40).} Under what conditions on \(f\) does \(S_\epsilon(f,P)\) have density \(1\) among the primes for every \(\epsilon>0\)?

The supplied source does not specify the meaning of density. Under the natural-density interpretation, the requested conclusion is
\[
\forall\epsilon>0,\qquad
\lim_{x\to\infty}
\frac{\#\{p\le x:p\in S_\epsilon(f,P)\}}{\pi(x)}=1,
\qquad \pi(x)=\#\{p\le x:p\text{ is prime}\}.
\]
This limit is an explicit interpretation, not a source-specified definition. Restricting to good reduction makes the reduced dynamics well defined; omitting finitely many primes does not affect natural or Dirichlet density. The source first fixes \(P\), but does not explicitly say whether the sought conditions on \(f\) must ensure the conclusion for every rational point with Zariski-dense orbit. No such additional universal quantifier is imposed here. The source also leaves the range of \(N\) implicit; the usual positive-dimensional interpretation is \(N\ge1\).
\subsection{Short English statement}
For a rational projective morphism of degree at least two and a rational point with Zariski-dense forward orbit, ask which conditions on the morphism ensure that, for every \(\epsilon>0\), the reduced orbit has at most \(p^{(1/2+\epsilon)N}\) distinct points for a density-one set of primes. The density convention and whether the conclusion must hold for every such starting point remain unspecified.
\subsection{Sources}
\begin{itemize}
\item Local UnsolvedMath registry, permanent ID \(800417\), exact source problem code \texttt{catalogue entry}; supplied complete source record, upstream revision \texttt{d371ee917ac9647fbb06cd6fc529deec6f3dda82}. The registry labels the problem \emph{open}; this is a supplied status, not an independently verified mathematical or literature assessment.
\item Benedetto, DeMarco, Ingram, Jones et al., \emph{Current Trends and Open Problems in Arithmetic Dynamics} (2018), Question 40. Supplied source URL: \url{https://arxiv.org/abs/1806.04980}. The record attributes the problem to these authors and gives proposed year 2018.
\item Supplied provenance: source accessed \texttt{2026-10-16}; extraction method \texttt{arxiv\_latex}. This access date and provenance are supplied metadata, not evidence of independent source consultation in preparing this draft.
\end{itemize}
\end{document}
